\scanpage{219}
\mspage{74}

ci-dessus (p. 69), de façon évidente : pour $A$ fixé, les $B$ tels que $(A,B)\in\overline R$ sont ceux qui sont contenus dans $A^{\unsure{\cup|\smile}}$, donc l'ensemble $\overline{\unsure{R|Q}}(A)$ n'est autre que $\overline Q_{/A^\cup}$, et dualement $\,{}^{\unsure{s|t}}\overline R(B)$ n'est autre que $\overline P_{/B^{\unsure{\cap|\frown}}}$. Ainsi

\[
 \beta(A)=A^\cup,\qquad\alpha(B)=B^\cap.
 \tag*{(5.14)}
\]

Si j'ai songé à introduire \deleted{\unsure{dans}} l'axiome $D_0$, pour \deleted{des} \inserted{un couple} d'ensembles ordonnés quelconques, c'est à cause des cas où ceux-ci ne sont \uline{pas} de la forme $\unsure{\mathfrak P|\mathcal P}(P_{\unsure{0|\sharp|,}}),\mathfrak P(Q_{\unsure{,|0}})$, et (alors \unsure{même|mêmes} qu'ils le seraient) où la relation $\unsure{R}$ \inserted{de base} n'est pas déduite \inserted{comme ci-dessus} d'une relation quelconque entre ensembles $P,Q$.

Je reviens \unsure{donc|sur} à l'axiomatique du début, avec $P,Q$ ordonnés, et $R$ satisfaisant (D). Dans tous les cas auxquels je peux songer, $P$ et $Q$ sont stables par Sup quelconques, donc (D) équivaut à $(D_0)$.

Notons que l'application

\[
 P\longrightarrow\mathfrak P(P),\qquad p\unsure{\vdash|\mapsto} P_{/p}
\]

est injective, et de $=$ pour $Q\to\mathfrak P(Q)$. Ainsi, la connaissance de $A\mapsto A^\cup$ équivaut à celle de $A\longmapsto\beta(A)$

\[
 \beta:\mathfrak P(P)\longrightarrow Q,
\]

et de $=$ celle de $B\vdash B^\cap$ équivaut à
% archive sheet number: 219
